\setcounter{section}{51}
\section{Sistemas de la Información}

Nombre completo: Sistemas de altas prestaciones. Grid Computing. Mainframe.

\localtableofcontents

\subsection{Tipos de sistemas de altas prestaciones}

También llamados High Performance Computing. Se pueden clasificar en varios tipos:
\begin{description}
    \item[Centralizado] Grandes servidores monolíticos de tipo SIMD o MIMD. Son difíciles de manejar y se crean de forma específica para la organización en ka que prestan servicio.
    \item[Cluster] Sistemas de distribución paralela y distribuida en un ámbito local.
    \item[Distribuido] Sistemas de distribución paralela en ámbito no localizado.
\end{description}

\subsubsection{Taxonomía de Flynn}
Es una clasificación de arquitecturas de computadores. Se basan en el número de instrucciones concurrentes y en los flujos de datos disponibles. 

Algunos ejemplos de estas clasificaciones son
\begin{description}
    \item[SISD] Von Neuman
    \item[SIMD] Gráficas y arquitectrura vectorial, minadores de bitcoins
    \item[MISD] Paralelismo redundante, navegación aérea, poco uso
    \item[MIMD] Sistemas distribuidos.
    \begin{description}
        \item[Memoria compartida] Se ejecuta un acceso simétrico a la memoria desde cualquier procesador
        \item[Memoria distribuida] Asignan una porción de memoria a los procesadores
        \item[Memoria distribuida compartida] Engloban los dos anteriores con memoria compartida y porciones distribuidas
        \item[Memoria solo caché] Procesador que solo usa caché como memoria, es un caso especial.
    \end{description}
\end{description}
\begin{table}[H]
\begin{tabular}{rcc}
                        & Una instrucción & Multiples instrucciones \\
    Un dato             & \texttt{SISD} & \texttt{MISD} \\
    Múltioples datos    & \texttt{SIMD} & \texttt{MIMD} \\
\end{tabular}
\end{table}

\subsection{Mainframes}
Sistemas multiusuario, multitarea y multiprocesador, orientado al procesamiento de grandes volúmenes de información y cálculo científico

\subsection{Sistemas con múltiples computadoras}
Son arquitecturas escalables que permiten un comportamiento dinámico.

\subsubsection{Clustering}
Arquitectura paralela distribuida formada por un conjunto de ordenadores independientes interconectados operando como un único recurso computacional. Se forma por nodos de computación, red de interconexión y un midleware que hace de interfaz proveiendo un único recurso computacional.

\subsubsection{Grid computing}
Sistema distribuido y paralelo que permite compartir, seccionar y agregar recursos de forma dinámica y en tiempo real. Estos recursos pueden asignarse dependiendo de disponibilidad, capacidad, rendimiento, o cualquier otro criterio.

El ejemplo mas famoso de Grid computing es el proyecto SETI.